\documentclass[12pt]{article}
\usepackage{amsmath}
\usepackage{amssymb}
\usepackage[margin=1in]{geometry}
\usepackage{graphicx}
\usepackage{hyperref}
\usepackage{xcolor}
\usepackage{booktabs}

% Placeholder macro for writer-fill fields
\newcommand{\fillme}[1]{\textcolor{red}{[\textit{#1}]}}

% Base path for all generated frames (frames7 run on cluster)
% Update SAMPLESNAP to whichever snapshot is representative (default: 0200)
\newcommand{\framesdir}{/scratch/vasissua/SHIVAN/analysis/frames7}
\newcommand{\samplesnap}{0200}

\title{Draft Figure Captions\\[0.5em]
\large FIRE+STARFORGE Disk Analysis Pipeline (m12f, $z\approx21$)}
\author{}
\date{}

\begin{document}
\maketitle
\tableofcontents
\newpage

%% ============================================================
\section{Master frames (per snapshot)}
%% ============================================================

These figures are produced for every snapshot by \texttt{notebooks/make\_disk\_movie\_frames.py}
via \texttt{ntbk.main(Defaults())} in \texttt{jupytertest3.py}.
The disk identification (center, orientation, membership) is recomputed independently
for each snapshot.

% --------------------------------------------------
\subsection{Frame A: surface density and velocity maps}
% --------------------------------------------------

\begin{figure}[h]
  \centering
  \includegraphics[width=\textwidth]{\framesdir/frame_maps_\samplesnap.png}
  \caption{
    \textbf{Spatial maps of the protostellar disk at snapshot \samplesnap.}
    All panels share a dark background and an inferno (surface density) or viridis
    (velocity) colour scale.
    The disk centre is placed at the most massive sink particle; the reference frame is
    rotated so that the disk angular momentum vector $\hat{L}$ points along $+z$
    (face-on view) or $+y$ (edge-on view).
    %
    \textit{Row 1 (surface density):}
    (a) Face-on projected surface density $\Sigma$ [M$_\odot\,\mathrm{pc}^{-2}$]
        over a $\sim600\,\mathrm{AU}$ box ($\log_{10}$, inferno);
    (b) same quantity over a $10\times$ wider field ($\sim6000\,\mathrm{AU}$)
        showing the surrounding envelope;
    (c) edge-on surface density (clean, no overlay);
    (d) edge-on surface density with disk membership overlay
        (particles satisfying the geometric+kinematic cuts highlighted).
    $\Sigma$ is computed with Meshoid kernel-weighted surface-density projection.
    %
    \textit{Row 2 (turbulent velocity):}
    (e) Face-on projection of the mass-weighted mean residual speed
        $|\delta v| = \sqrt{\delta v_r^2 + \delta v_\phi^2 + \delta v_z^2}$\,[km\,s$^{-1}$]
        after subtracting 20-bin mass-weighted radial streaming profiles of
        $\langle v_r\rangle(r)$ and $\langle v_\phi\rangle(r)$;
    (f) face-on projection of the mass-weighted standard deviation
        $\sigma_{|\delta v|}$\,[km\,s$^{-1}$] of the same quantity;
    (g,\,h) edge-on equivalents of panels (e,\,f).
    Colour scale is linear (viridis).
    %
    \textit{Row 3 (midplane and magnetic field):}
    (i) Face-on surface density of the midplane slice: only particles with
        $|z_\mathrm{fo}| < \max(\tilde{h}, P_5(|z_\mathrm{fo}|))$
        are included, where $\tilde{h}$ is the median smoothing length
        and $P_5$ is the 5th percentile of $|z_\mathrm{fo}|$;
    (j) face-on mass-weighted $|B_z|$ map [Gauss] (averaged along the line of
        sight using the Meshoid kernel; requires \texttt{MagneticField} in snapshot);
    (k,\,l) reserved for future use.
    %
    \fillme{Describe the overall morphology, any spiral arms, fragmentation,
            or asymmetries visible in the surface density maps.}
    \fillme{Comment on the spatial extent of the turbulent velocity field
            relative to the identified disk boundary.}
    \fillme{Compare the midplane slice (i) to the projected map (a) --
            does the midplane show a higher central concentration?}
  }
  \label{fig:frame_maps}
\end{figure}

% --------------------------------------------------
\subsection{Frame B: kinematic scatter plots and radial profiles}
% --------------------------------------------------

\begin{figure}[h]
  \centering
  \includegraphics[width=\textwidth]{\framesdir/frame_analysis_\samplesnap.png}
  \caption{
    \textbf{Kinematic analysis panels for snapshot \samplesnap.}
    %
    \textit{Row 1 (particle-level scatter):}
    (a) Radial velocity $v_r$\,[km\,s$^{-1}$] vs.\ cylindrical radius $r$\,[AU]
        for every disk particle (scatter), with the 20-bin mass-weighted mean
        $\langle v_r\rangle(r)$ overlaid in red;
    (b) azimuthal velocity $v_\phi$\,[km\,s$^{-1}$] vs.\ $r$ with
        $\langle v_\phi\rangle(r)$ (red) and the Keplerian profile
        $v_K = \sqrt{G M_\mathrm{enc}(<r)/r}$ (yellow dashed);
    (c) residual speed $|\delta v|$\,[km\,s$^{-1}$] vs.\ $r$
        (streaming-subtracted, 20 bins) with profile;
    (d) Toomre-$Q$ map (face-on, colour scale centred at $Q=1$).
    %
    \textit{Row 2 (annular profiles):}
    (e) Radially averaged Toomre $Q(r)$ with the $Q=1$ critical line (dashed),
        where $Q = \sigma_r \kappa / (\pi G \Sigma)$
        and the epicyclic frequency $\kappa \approx \Omega_K$;
    (f) surface density profile $\Sigma(r)$\,[M$_\odot\,\mathrm{pc}^{-2}$] (log scale);
    (g) sound-speed $c_s$, radial velocity dispersion $\sigma_r$, and
        mean residual speed $\langle|\delta v|\rangle$ as functions of $r$ [km\,s$^{-1}$];
    (h) Mach number $\mathcal{M} = |\delta v|/c_s$ as a function of $r$.
    %
    \textit{Row 3 (stability diagnostics):}
    (i) Mass-weighted mean resolution $\Delta x = (m/\rho)^{1/3}$ [AU] per annulus
        (log scale) with the Jeans length $\lambda_J$ overlaid;
    (j) Virial parameter
        $\alpha(r) = v_\mathrm{turb}^2\,r\,/\,(G\,M_\mathrm{enc}(<r))$
        per annulus, where $v_\mathrm{turb}^2$ is the mass-weighted mean squared
        residual speed and $M_\mathrm{enc}$ includes both sink and gas mass;
    (k) dimensionless mass-to-flux ratio
        $\mu(r) = 2\pi\sqrt{G}\,\Sigma / |B_z|$ per annulus
        (critical value $\mu_\mathrm{crit}=1$ shown as dashed line;
        $\mu>1$ magnetically supercritical);
    (l) reserved for future use.
    %
    \fillme{State whether $v_\phi(r)$ closely follows the Keplerian profile,
            indicating a rotationally supported disk.}
    \fillme{Report the radial range over which $Q < 1$, if any, and
            whether this correlates with the locations of sink particles.}
    \fillme{Comment on whether the disk is resolved at all radii shown
            (compare $\Delta x$ to $r$) and where resolution limits apply.}
    \fillme{Discuss whether $\alpha < 1$ (sub-virial, still contracting)
            or $\alpha \sim 1$ (virialized) across the disk.}
    \fillme{Discuss whether $\mu > 1$ (magnetically supercritical) throughout
            or whether there is a magnetically subcritical core.}
  }
  \label{fig:frame_analysis}
\end{figure}

% --------------------------------------------------
\subsection{Frame C: phase diagrams (optional)}
% --------------------------------------------------

\begin{figure}[h]
  \centering
  % Frame C is only produced when include_phase_in_master=True in Defaults()
  % \includegraphics[width=\textwidth]{\framesdir/frame_phase_\samplesnap.png}
  \fbox{\parbox{0.9\textwidth}{\centering
    \textit{Phase diagram frame (frame\_phase\_\samplesnap.png).
    Produced only when \texttt{Defaults.include\_phase\_in\_master = True}.}}}
  \caption{
    \textbf{Thermochemical phase diagrams for snapshot \samplesnap.}
    (a) Temperature $T$\,[K] vs.\ gas density $n$\,[cm$^{-3}$] for all disk
        particles, colour-coded by $H_2$ fraction $f_{H_2}$;
    (b) $\log_{10} f_{H_2}$ vs.\ $n$ for the same particles.
    Both quantities are derived from the \texttt{InternalEnergy} and
    \texttt{Density} fields; molecular fraction is tracked by the
    STARFORGE chemistry network.
    Dashed lines indicate fiducial temperature thresholds ($T=100\,\mathrm{K}$,
    $T=10\,\mathrm{K}$) and density thresholds used in the disk identification.
    %
    \fillme{Describe the dominant thermal phase(s) of the disk gas
            (warm neutral, cold neutral, molecular).}
    \fillme{Report the typical $T$--$n$ locus of the identified disk particles
            and compare to the surrounding envelope.}
  }
  \label{fig:frame_phase}
\end{figure}

%% ============================================================
\section{Individual / thematic frames (per snapshot)}
%% ============================================================

These figures are produced alongside the master frames when \texttt{data\_outdir}
is set (always the case in \texttt{jupytertest3.py}).
They are saved to \texttt{frames7/individual\_frames/} and are intended
to be inserted individually throughout a paper.

% --------------------------------------------------
\subsection{Density}
% --------------------------------------------------

\begin{figure}[h]
  \centering
  \includegraphics[width=\textwidth]{\framesdir/individual_frames/frame_density_\samplesnap.png}
  \caption{
    \textbf{Surface density views at snapshot \samplesnap.}
    Five-panel figure showing projected surface density
    $\Sigma$\,[M$_\odot\,\mathrm{pc}^{-2}$] (inferno colour scale, logarithmic).
    From left to right:
    (a) face-on, $\sim600\,\mathrm{AU}$ box (disk scale);
    (b) face-on, $\sim6000\,\mathrm{AU}$ box (envelope scale);
    (c) edge-on, clean projection;
    (d) edge-on with disk membership overlay;
    (e) face-on midplane slice ($|z_\mathrm{fo}| < \tilde{h}$).
    Sink particles are marked with white $\star$ symbols scaled to their mass.
    Surface density is computed using Meshoid kernel-weighted projection at
    resolution $400\times400$ pixels over the specified field of view.
    %
    \fillme{Describe the disk aspect ratio inferred from the edge-on view.}
    \fillme{Note whether the envelope (panel b) shows infalling streamers
            or tidal distortions.}
    \fillme{Comment on the central concentration visible in the midplane
            slice (panel e) relative to the column-integrated view (panel a).}
  }
  \label{fig:density}
\end{figure}

% --------------------------------------------------
\subsection{Velocity dispersion maps}
% --------------------------------------------------

\begin{figure}[h]
  \centering
  \includegraphics[width=\textwidth]{\framesdir/individual_frames/frame_velocity_\samplesnap.png}
  \caption{
    \textbf{Turbulent velocity maps at snapshot \samplesnap.}
    Four-panel figure (viridis colour scale, linear).
    (a) Face-on projection of the mass-weighted mean residual speed
        $|\delta v|$\,[km\,s$^{-1}$];
    (b) face-on projection of the standard deviation
        $\sigma_{|\delta v|}$\,[km\,s$^{-1}$];
    (c,d) edge-on equivalents of (a,b).
    Streaming motions are subtracted by computing 20-bin mass-weighted radial
    profiles of $v_r(r)$ and $v_\phi(r)$ within the disk plane, so
    $\delta v_r = v_r - \langle v_r\rangle_\mathrm{bin}$,
    $\delta v_\phi = v_\phi - \langle v_\phi\rangle_\mathrm{bin}$,
    $\delta v_z = v_z$,
    and $|\delta v| = (\delta v_r^2 + \delta v_\phi^2 + \delta v_z^2)^{1/2}$.
    %
    \fillme{Describe whether the turbulence is centrally concentrated or
            extends to large radii.}
    \fillme{Compare the face-on vs.\ edge-on velocity dispersion structures
            -- is there a systematic vertical gradient?}
  }
  \label{fig:velocity}
\end{figure}

% --------------------------------------------------
\subsection{Kinematics profiles}
% --------------------------------------------------

\begin{figure}[h]
  \centering
  \includegraphics[width=\textwidth]{\framesdir/individual_frames/frame_kinematics_\samplesnap.png}
  \caption{
    \textbf{Radial kinematic profiles at snapshot \samplesnap.}
    Three-panel scatter+profile plot for all disk-identified gas particles.
    (a) Radial velocity $v_r$\,[km\,s$^{-1}$] vs.\ cylindrical radius $r$\,[AU]:
        grey scatter is individual particles; red curve is the 20-bin
        mass-weighted mean $\langle v_r\rangle(r)$;
    (b) azimuthal velocity $v_\phi$ vs.\ $r$, with $\langle v_\phi\rangle(r)$ (red)
        and Keplerian profile $v_K = \sqrt{G M_\mathrm{enc}/r}$ (yellow dashed);
    (c) residual speed $|\delta v|$ vs.\ $r$ after streaming subtraction,
        with the 20-bin mass-weighted mean profile overlaid.
    In all panels the enclosed mass $M_\mathrm{enc}(<r) = M_* + M_\mathrm{gas,cum}$
    includes both sink particle(s) and cumulative disk gas mass.
    %
    \fillme{State whether $v_\phi(r)$ is consistent with Keplerian rotation
            and over what radial range.}
    \fillme{Report the sign and magnitude of $\langle v_r\rangle(r)$:
            negative (infall) or positive (outflow/expanding)?}
    \fillme{Describe the radial profile of $|\delta v|$: does it peak
            near the centre or at an intermediate radius?}
  }
  \label{fig:kinematics}
\end{figure}

% --------------------------------------------------
\subsection{Toomre stability}
% --------------------------------------------------

\begin{figure}[h]
  \centering
  \includegraphics[width=\textwidth]{\framesdir/individual_frames/frame_toomre_\samplesnap.png}
  \caption{
    \textbf{Toomre gravitational stability at snapshot \samplesnap.}
    (a) Face-on map of the local Toomre $Q$ parameter:
    \begin{equation}
      Q(r) = \frac{\sigma_r\,\kappa}{\pi G \Sigma},
    \end{equation}
    where $\sigma_r$ is the radial velocity dispersion, $\kappa \approx \Omega_K$
    is the epicyclic frequency (Keplerian approximation), and $\Sigma$ is the
    local surface density.
    Regions with $Q < 1$ are gravitationally unstable to fragmentation
    (colour scale diverging at $Q=1$).
    (b) Annularly averaged $Q(r)$ profile with the critical line $Q=1$ (dashed).
    The radial bins are linearly spaced from $0$ to the 95th-percentile radius of
    the disk; $Q$ is computed from bin-averaged $\sigma_r$, $\Omega$, and $\Sigma$.
    Sink particle positions are shown as white stars in panel (a).
    %
    \fillme{Report the minimum value of $Q$ and its location.}
    \fillme{Describe whether gravitationally unstable regions ($Q < 1$)
            coincide spatially with existing or newly forming sink particles.}
    \fillme{Note whether $Q$ increases monotonically with $r$ or shows
            a local minimum at intermediate radii (ring instability).}
  }
  \label{fig:toomre}
\end{figure}

% --------------------------------------------------
\subsection{Radial profiles}
% --------------------------------------------------

\begin{figure}[h]
  \centering
  \includegraphics[width=\textwidth]{\framesdir/individual_frames/frame_profiles_\samplesnap.png}
  \caption{
    \textbf{Annular radial profiles at snapshot \samplesnap.}
    (a) Volume-averaged density $\rho(r)$\,[g\,cm$^{-3}$] (log scale) as a
        function of cylindrical radius $r$\,[AU], for disk-identified particles
        binned into 20 equal-width annuli;
    (b) velocity scales [km\,s$^{-1}$]: sound speed
        $c_s = \sqrt{\gamma P/\rho}$ with $\gamma = 5/3$,
        radial velocity dispersion $\sigma_r$,
        and mean residual speed $\langle|\delta v|\rangle$;
    (c) sonic Mach number $\mathcal{M} = \langle|\delta v|\rangle/c_s$
        (log scale; $\mathcal{M} = 1$ dashed);
    (d) mass-weighted mean spatial resolution
        $\Delta x = (m/\rho)^{1/3}$\,[AU] per annulus (log scale),
        where $m$ is the particle mass and $\rho$ the local density.
    %
    \fillme{Describe the density power-law slope and whether it is
            consistent with a self-similar collapse profile ($\rho \propto r^{-3/2}$
            for free-fall, or $\rho \propto r^{-1}$ for a Mestel disk).}
    \fillme{Report whether the disk is sub-sonic ($\mathcal{M} < 1$) or
            super-sonic ($\mathcal{M} > 1$) and at what radii.}
    \fillme{State the radius at which the resolution $\Delta x$ begins
            to exceed the local Jeans length or scale height, limiting
            the fidelity of the simulation.}
  }
  \label{fig:profiles}
\end{figure}

% --------------------------------------------------
\subsection{Stability diagnostics}
% --------------------------------------------------

\begin{figure}[h]
  \centering
  \includegraphics[width=\textwidth]{\framesdir/individual_frames/frame_stability_\samplesnap.png}
  \caption{
    \textbf{Dynamical stability diagnostics at snapshot \samplesnap.}
    (a) Virial parameter profile:
    \begin{equation}
      \alpha(r) = \frac{v_\mathrm{turb}^2\,r}{G\,M_\mathrm{enc}(<r)},
    \end{equation}
    where $v_\mathrm{turb}^2$ is the mass-weighted mean squared residual velocity
    per annulus and $M_\mathrm{enc}$ includes sink mass plus enclosed disk gas.
    The horizontal dashed line marks $\alpha = 1$ (virial equilibrium);
    $\alpha < 1$ indicates sub-virial (still contracting) gas;
    $\alpha > 2$ indicates unbound gas.
    (b) Dimensionless mass-to-flux ratio profile:
    \begin{equation}
      \mu(r) = \frac{2\pi\sqrt{G}\,\Sigma}{|B_z|},
    \end{equation}
    where $\Sigma$ is the annulus-averaged column density [g\,cm$^{-2}$]
    and $|B_z|$ is the mass-weighted mean absolute value of the face-on
    magnetic field component [Gauss].
    The critical value $\mu_\mathrm{crit} = 1$ (magnetically critical) is shown
    as a dashed line; $\mu > 1$ is magnetically supercritical (gravity dominates
    over magnetic support).
    %
    \fillme{Describe the virial state of the disk: is $\alpha$ roughly uniform
            or does it show a strong radial gradient?}
    \fillme{Report whether the disk is globally magnetically supercritical
            ($\mu > 1$ everywhere) or whether a magnetically subcritical region
            exists and at what radius.}
  }
  \label{fig:stability}
\end{figure}

% --------------------------------------------------
\subsection{Magnetic field (conditional)}
% --------------------------------------------------

\begin{figure}[h]
  \centering
  % Only produced if MagneticField is present in the snapshot
  % \includegraphics[width=\textwidth]{\framesdir/individual_frames/frame_bfield_\samplesnap.png}
  \fbox{\parbox{0.9\textwidth}{\centering
    \textit{Magnetic field frame (frame\_bfield\_\samplesnap.png).
    Produced only when \texttt{MagneticField} is present in the loaded snapshot fields.}}}
  \caption{
    \textbf{Magnetic field structure at snapshot \samplesnap.}
    (a) Face-on map of the mass-weighted mean $|B_z|$ [Gauss],
        projected along the disk normal axis using Meshoid
        kernel-weighted averaging;
    (b) radial profile of the mass-to-flux ratio $\mu(r)$
        (same quantity as in Figure~\ref{fig:stability}b, shown here for
        direct comparison with the $|B_z|$ map).
    %
    \fillme{Describe the spatial structure of $B_z$: is it centrally
            concentrated, does it show a coherent large-scale field, or
            is it disordered?}
    \fillme{Compare the magnetic field morphology to the surface density
            map (Figure~\ref{fig:density}a) -- does $|B_z|$ peak where
            the density is highest?}
  }
  \label{fig:bfield}
\end{figure}

%% ============================================================
\section{Evolution plots (one figure per pipeline run)}
%% ============================================================

These figures are produced once per \texttt{jupytertest3.py} run,
covering all available snapshots.

% --------------------------------------------------
\subsection{Mass evolution}
% --------------------------------------------------

\begin{figure}[h]
  \centering
  \includegraphics[width=0.75\textwidth]{\framesdir/mass_evolution.png}
  \caption{
    \textbf{Mass evolution of the protostellar system.}
    Time axis is $t - t_1$\,[kyr], where $t_1$ is the formation time of
    the first sink particle.
    (a) Stellar mass $M_*$ (sum of all sink masses, yellow), disk gas mass
        $M_\mathrm{gas,disk}$ from the geometric+kinematic disk identification
        algorithm (cyan solid), and gas mass within a fixed spherical aperture of
        radius $5 r_\mathrm{max}$ (cyan dashed).
        A power-law fit $M_* \propto t^\alpha$ is overlaid in red with the
        best-fit exponent $\alpha$; the reference $M_* \propto t^2$ curve
        (expected for spherical Bondi-like accretion) is shown dotted.
    (b) Star formation efficiency $f_* = M_* / (M_* + M_\mathrm{gas})$
        for both disk-identification (solid) and fixed-aperture (dashed) denominators.
    (c) Stellar accretion rate $\dot{M}_*$\,[M$_\odot$\,yr$^{-1}$] (solid) and
        gas depletion rate $-\dot{M}_\mathrm{gas,apert}$\,[M$_\odot$\,yr$^{-1}$] (dashed).
    (d) Sphere-of-influence radius $r_\mathrm{SOI}$\,[AU], defined as the radius
        at which the enclosed gas mass equals $M_*$.
    %
    \fillme{Report the power-law index $\alpha$ of the stellar mass growth and
            compare to theoretical expectations.}
    \fillme{Describe the star formation efficiency at the end of the simulation
            and whether it is still rising.}
    \fillme{Note whether the accretion rate is approximately constant or
            shows episodic bursts.}
    \fillme{Comment on the evolution of $r_\mathrm{SOI}$ -- does the sink's
            gravitational sphere of influence grow monotonically?}
  }
  \label{fig:mass_evolution}
\end{figure}

% --------------------------------------------------
\subsection{Energy evolution}
% --------------------------------------------------

\begin{figure}[h]
  \centering
  \includegraphics[width=0.75\textwidth]{\framesdir/energy_evolution.png}
  \caption{
    \textbf{Mechanical energy evolution of disk gas.}
    Energies are computed for disk-identified particles only;
    time axis is $t - t_1$\,[kyr] as in Figure~\ref{fig:mass_evolution}.
    Units are $10^{44}$\,erg.
    (a) Energy components (log scale):
        rotational kinetic energy
        $E_\mathrm{rot} = \tfrac{1}{2}\sum_i m_i v_{\phi,i}^2$,
        turbulent kinetic energy
        $E_\mathrm{turb} = \tfrac{1}{2}\sum_i m_i |\delta v_i|^2$,
        total kinetic energy
        $E_\mathrm{kin} = E_\mathrm{rot} + E_\mathrm{turb}$,
        and gravitational well depth $|E_\mathrm{pot}|$ (tomato line).
        The potential energy is computed using an annular-binned cylindrical
        approximation:
        $E_\mathrm{pot} \approx -G \sum_b M_\mathrm{enc}(<r_b)\,M_b / r_b$,
        where the sum is over 20 radial bins with $r_b$ the bin-centre radius
        (softened to at least the mean smoothing length to avoid singularities).
    (b) Magnitude of total mechanical energy
        $|E_\mathrm{tot}| = |E_\mathrm{kin} + E_\mathrm{pot}|$.
        Because $|E_\mathrm{pot}| \gg E_\mathrm{kin}$ for a bound, sub-virial disk,
        $E_\mathrm{tot} < 0$ (solid curve shows $|E_\mathrm{tot}|$;
        a dashed style would indicate an unbound system).
    (c) Virial ratio $2E_\mathrm{kin}/|E_\mathrm{pot}|$:
        $= 1$ at virial equilibrium (red dashed); $< 1$ for a contracting disk;
        $= 2$ for an unbound system (orange dotted).
    %
    \fillme{Report whether the disk is sub-virial ($2E_k/|E_p| < 1$)
            throughout its evolution and whether it approaches virial equilibrium.}
    \fillme{Describe whether $E_\mathrm{rot}$ and $E_\mathrm{turb}$ grow at
            the same rate (suggesting similar channels of energy injection) or
            differently.}
    \fillme{Comment on whether $|E_\mathrm{pot}|$ growing faster than
            $E_\mathrm{kin}$ is consistent with ongoing gravitational collapse.}
  }
  \label{fig:energy_evolution}
\end{figure}

% --------------------------------------------------
\subsection{Toomre-$Q$ heatmap}
% --------------------------------------------------

\begin{figure}[h]
  \centering
  \includegraphics[width=0.75\textwidth]{\framesdir/Q_heatmap.png}
  \caption{
    \textbf{Space--time Toomre $Q$ heatmap.}
    Colour encodes the annularly-averaged Toomre $Q(r,t)$
    as a function of cylindrical radius $r$\,[AU] (vertical axis) and
    time $t - t_1$\,[kyr] (horizontal axis).
    Profiles are assembled from all snapshots for which disk-classified particles
    were found; each profile is saved to \texttt{qprofiles/qprofile\_XXXX.npz}
    during frame rendering and then co-assembled by \texttt{make\_Q\_heatmap()}.
    The critical contour $Q = 1$ (gravitational instability threshold) is
    drawn in white.
    %
    \fillme{Identify and describe the epoch(s) at which $Q < 1$ regions appear
            and whether they propagate inward or outward over time.}
    \fillme{Comment on whether the gravitationally unstable zone persists or
            is stabilised by accretion heating.}
  }
  \label{fig:Q_heatmap}
\end{figure}

%% ============================================================
\section{Supplementary per-snapshot plots}
%% ============================================================

% --------------------------------------------------
\subsection{Phase diagrams}
% --------------------------------------------------

\begin{figure}[h]
  \centering
  \includegraphics[width=\textwidth]{\framesdir/T_H2_rho_phase_plots/phase_\samplesnap.png}
  \caption{
    \textbf{Thermochemical phase diagram at snapshot \samplesnap.}
    Scatter plot of gas temperature $T$\,[K] vs.\ number density
    $n$\,[cm$^{-3}$] (both axes logarithmic) for all gas particles
    in the cutout snapshot, colour-coded by molecular hydrogen fraction
    $f_{H_2}$ (viridis).
    Temperature is derived from the GIZMO internal energy field as
    $T = (\gamma - 1)\,\mu\,m_H\,u / k_B$, where $u$ is the specific
    internal energy, $\mu$ is the mean molecular weight, and $\gamma = 5/3$.
    The $H_2$ fraction is tracked by the FIRE-3/STARFORGE chemistry network
    and stored in the snapshot.
    Disk-identified particles (satisfying the geometric + kinematic cuts)
    are highlighted in a contrasting colour.
    %
    \fillme{Describe the dominant thermal phase of the disk gas and whether
            it is predominantly molecular ($f_{H_2} > 0.5$).}
    \fillme{Identify any warm ($T > 100$\,K) dense gas that may indicate
            accretion shocks or proto-stellar irradiation.}
  }
  \label{fig:phase}
\end{figure}

% --------------------------------------------------
\subsection{Velocity power spectrum}
% --------------------------------------------------

\begin{figure}[h]
  \centering
  \includegraphics[width=\textwidth]{\framesdir/velocity_power_spectra/plots/vps_\samplesnap.png}
  \caption{
    \textbf{Velocity power spectrum at snapshot \samplesnap.}
    Power spectral density $P(k)$ of the three-dimensional velocity field of
    disk-identified gas, computed via a non-uniform fast Fourier transform
    (NFFT) or interpolation to a regular grid, as a function of wavenumber
    $k$\,[$\mathrm{AU}^{-1}$].
    The Kolmogorov slope $P \propto k^{-5/3}$ and the Burgers (supersonic)
    slope $P \propto k^{-2}$ are overlaid for reference.
    The vertical dashed line marks the mean particle resolution scale
    $k_\mathrm{res} = 1/\Delta x$.
    %
    \fillme{Report the inertial-range slope of the power spectrum and
            whether it is closer to Kolmogorov or Burgers.}
    \fillme{Identify the injection scale (turnover to shallower slope at
            low $k$) and compare it to the disk radius.}
    \fillme{Note the dissipation scale relative to the numerical resolution
            limit $k_\mathrm{res}$.}
  }
  \label{fig:vps}
\end{figure}

\end{document}
